\part{底层认知：ADHD 是一套算法，不是一处故障}

\section{ADHD 是一套算法，不是一处故障}

先聊命名。\term{「注意力缺陷多动障碍」}这个译名，把所有人都带偏了——它让你以为 ADHD 就是「注意力不够、太多动」。于是 ADHD 本人也天天自责：「我要是能再专注一点就好了。」但真正的图景，和「专注力不够」几乎没有关系。

原视频给了一个更好的框架，我们称它为\term{算法观}（BV13HoGYiEgF）：

把大脑理解成一套算法。普通大脑的成本结构是：深度运算（单点纵深、线性推进）便宜，主动注意力便宜；而 ADHD 大脑恰好反过来——\term{广度运算便宜，深度运算贵；被动注意力便宜，主动注意力贵}。

这不是「坏了」，是「成本结构不同」。就像 GPU 和 CPU：GPU 跑并行便宜，跑单线程反而贵；CPU 正好相反。你不能指着 GPU 说「你怎么这么不会跑单线程」，它是另一种东西。

顺着这个框架，很多现象第一次能解释清楚了：

\begin{itemize}
  \item \term{为什么能刷六个小时视频，却不能看十分钟书}——刷视频用的是便宜的被动注意力，看书用的是昂贵的主动注意力。
  \item \term{为什么上课听不懂就彻底崩了}——课堂要求连续线性跟随，那是深度运算，是 ADHD 最贵的项目。
  \item \term{为什么「重要」的事情推不动}——「重要」是理性的语言，而你的启动器只听多巴胺的语言。
\end{itemize}

\subsection{椭圆竞争与超专注}

注意力空间里，无意识的念头（原视频叫「椭圆」）一直在互相竞争。普通人的主动注意力圈大，能把冒出来的念头按住；ADHD 的主动圈小，念头经常赢，这就是走神。

但当某一个念头强到占满几乎整个注意力空间，就会出现\term{超专注}——连续几小时沉浸，外界全部屏蔽。注意，这里的关键不是「能不能专注」，而是\term{电源开关不在你手上}：你没法在「该专注的时候」打开它，只能等它自己来。

超专注常见于游戏，但它可以被引导。原视频分享过一段真实经历：大三上学期，他能一天学十三个小时，是因为他当时真的相信「把定义和定理一个个弄清楚，一步步推，就能考进前几名」——\term{他看到了真凉}（实利）。「想考好成绩」这个念头足够强，于是在注意力竞争里赢了，连续数月霸占整个空间。

\begin{insight}
ADHD 自救的核心，从来不是「把主动注意力练大」（那是逼 GPU 去跑 CPU 的活，又贵又伤），而是\term{让被动注意力更多地落在正事上}——把那个免费、便宜、量大管饱的通道，引导到对的地方。原视频叫它\term{以广促深}。
\end{insight}

\subsection{核心特征与伴随特征：别把所有毛病都往 ADHD 上贴}

网上关于 ADHD 的描述，九成在讲坏处，而且讲得五花八门：记性差、粗心、忘时间、情绪化、拖延……看起来每一条都像「你」。原视频把这里叫\term{归并效应}（BV1zU9LBVEnm）：一个标签被广泛使用之后，那些标签并不必然导致的东西，也会被贴到这个标签上。

但真实情况是：\term{ADHD 的核心特征只有一样——执行功能缺陷}（启动困难、维持线性集中困难）。其余全是伴随特征，个体差异极大，甚至互相矛盾：有人记性差得要命，有人对感兴趣领域的知识记得极牢；有人情绪化，有人反而理性得过分。

所以「我是不是真的 ADHD」这个问题，本身可能问错了。更值得问的是：\term{这套算法，在我的生活里造成了多少实际损害？}按 DSM 的标准，确诊 = 特征 + 功能损害，两者缺一不可。同一个人，在清华当数学博士时每天写数学写得好好的，没法确诊；在博士期间连最简单的任务都启动不了、每天五点就想睡觉时，一测一个准。\term{ADHD 是一个「环境交互」的标签，不是一个贴在脑门上的本质属性}。

\begin{notehot}
这一段不是让你别去了解自己，恰恰相反——是让你别被「网上的 ADHD 画像」吓到或框死。你可以有 ADHD 的算法，同时记性极好；你可以没确诊，但这套方法对你也照样有效。标签只是地图，不是领地。
\end{notehot}

\key{你不是不能专注，你是「主动专注」这项服务对你来说太贵了。与其硬充，不如学会用便宜的通道。}
